\documentclass[11pt,a4paper]{article}
\usepackage[T1]{fontenc}
\usepackage[utf8]{inputenc}
\usepackage{lmodern,microtype,amsmath,amssymb,amsthm,booktabs,longtable,array}
\usepackage[margin=25mm]{geometry}
\usepackage{xcolor}
\usepackage{fancyhdr}
\usepackage{enumitem}
\usepackage{xurl}
\usepackage[colorlinks=true,linkcolor=blue!45!black,urlcolor=blue!45!black,citecolor=blue!45!black]{hyperref}
\usepackage{bookmark}
\setlength{\parindent}{0pt}
\setlength{\parskip}{5pt plus1pt}
\setlength{\emergencystretch}{3em}
\setlist{nosep}
\pagestyle{fancy}\fancyhf{}
\fancyhead[L]{\small Rank 11, multiplicity one}
\fancyhead[R]{\small Categorification audit}
\fancyfoot[C]{\thepage}
\setlength{\headheight}{14pt}
\newtheorem{theorem}{Theorem}
\newtheorem{proposition}{Proposition}
\newtheorem{lemma}{Lemma}
\providecommand{\FPdim}{\operatorname{FPdim}}
\providecommand{\Rep}{\operatorname{Rep}}
\providecommand{\Vec}{\operatorname{Vec}}
\providecommand{\BrPic}{\operatorname{BrPic}}
\providecommand{\Irr}{\operatorname{Irr}}
\providecommand{\Hom}{\operatorname{Hom}}
\providecommand{\End}{\operatorname{End}}
\providecommand{\id}{\operatorname{id}}
\providecommand{\Aut}{\operatorname{Aut}}
\providecommand{\rk}{\operatorname{rank}}
\providecommand{\tightlist}{\setlength{\itemsep}{0pt}\setlength{\parskip}{0pt}}
\begin{document}
\begin{center}
{\Large\bfseries Rank-eleven multiplicity-one fusion rings}\\[5pt]
{\large Exclusions, coherent models, and exact verification}\\[6pt]
23 October 2026
\end{center}
\section{Result and mathematical scope}
The input consists of 182 physical multiplication-table lines. All 182 satisfy
the fusion-ring axioms, and exact invariant refinement separates every pair.
There are 153 commutative and 29 noncommutative rings. Identifiers
\texttt{R11\_1\_i} always retain physical input line $i$; no sorting or basis
change is allowed to redefine an identifier.

\begin{center}\begin{tabular}{lr}\toprule Status & Rings\\\midrule

Excluded over $\mathbb C$ & 155\\

Unitary realization verified & 27\\

Unresolved after recorded attempts & 0\\

\bottomrule\end{tabular}\end{center}
An exclusion proves that no complex fusion category has the indicated based
ring. A model proves existence of a unitary category, including coherence of
the construction and an exact identification of all $11^3=1331$ coefficients.
Neither assertion classifies all possible associators of a realized ring.
No timeout, unsuccessful model search, floating-point test, or failed
relaxation is used to assign an exclusion.

\medskip
\noindent\textbf{AI-generated research \textperiodcentered\ independent verification pending.}
The proofs, code and data were developed with AI assistance. Independent
human verification and Lean formalization remain pending. Exact automated
checks establish only their stated scope.

\begin{theorem}[Classification of the supplied multiplicity-one census]
Of the 182 supplied rank-eleven multiplicity-one fusion rings, exactly 27
admit a fusion categorification over $\mathbb C$. Each admits a unitary
realization. The remaining 155 admit no complex fusion categorification.
\end{theorem}
The theorem concerns based-ring isomorphism classes in the supplied dataset.
Enumeration completeness is producer-reported: every duality stratum
completed and the complete source is supplied, while independent integration
checks establish tensor validity and uniqueness. The full enumeration was
not independently repeated. The census provenance is retained separately in
\url{https://github.com/sebastienpalcoux/fusion-ring-census}.
The canonical manuscript and per-input classification evidence are in
\url{https://github.com/sebastienpalcoux/grothendieck-ring-classification}.

\section{Known criteria and the order of the audit}
The audit first reused the literature and accepted lower-rank project methods:
formal-codegree d-numbers, character-trace cyclotomicity, primitive codegrees,
integral center dimensions, all fusion-subring indices, near-group and
Tambara--Yamagami restrictions, zero/one-spectrum obstructions, cyclic
isotypic integrality, module matrices, and inherited solvability arguments.
It then applied complete signed-sphericalization and scalar-character checks
before the center-dimension and reciprocity refinements. Model searches used
established families before constructing additional coherent actions.

The detailed criterion-by-criterion inventory is
\path{work/rank11/criteria_audit/CRITERIA_AUDIT.md}. It records hypotheses,
sources, implementation boundaries, and conditions which are only unitary or
braided. It covers the surveyed literature and the accepted project corpus;
it does not claim that every conceivable theorem or consequence of the
literature has been algorithmically decided.

The exact arithmetic audit checks 919 proper nontrivial fusion subrings,
286 distinct irreducible field polynomials, and the irreducible algebra-block
decomposition of all 29 noncommutative inputs. Cyclotomicity is checked on
character traces in the noncommutative case. The four arithmetic receipts give
97 excluded inputs. The cyclic test gives 66 exclusions with overlap. The
zero/one-spectrum and module proofs use actual algebra and pentagon
constraints, with independent integer replays. Overlapping proofs are retained
in the ledger and are counted only once.

The signed analysis covers all 59 nonintegral inputs surviving the initial
arithmetic and cyclic tests. It checks every admissible sign branch and every
possible spherical dimension character. In particular, the six mixed
dimension-40 branches are handled directly; they are not assumed
pseudo-unitary. The dimension-lattice and two-character arguments below use
that reduction as an explicit proof dependency.

For the final integral pair, all primary positivity criteria and the
noncommutative Schur-product criterion pass by exact matrix identities.
Nonnegative rational coupling certificates establish all of the tested
Minkowski inequalities for every real $p\geq1$ and every nonnegative input
vector. These are genuine all-parameter passes, not numerical searches for
witnesses. The new graded-extension argument below resolves the pair.
The independent Titsworth--Hagge data crosswalk also checks their rank-eleven
dimension-32 tensor: it is exactly input 30, already realized here.

\subsection{Evidence and reproducibility}
\path{work/rank11/ledger.json} is the per-input record of status, original
tensor hashes, proof files, certificates, and replay commands.
\path{work/rank11/ledger.csv} is its compact human-readable counterpart.
The input hash and mathematical assertions are bound to the original tensor,
not to approximate dimensions or an isomorphism signature alone.
Run the portable replay from the extracted package root:
\begin{quote}\ttfamily
python work/rank11/replay\_all.py
\end{quote}
The accompanying README documents supported options, dependencies and the
recorded complete replay. Discovery programs are kept distinct from accepted
certificate checks. Checks reject mismatches explicitly even under
\texttt{python -O}. The model audit independently reconstructs the highest-risk
coherence equations and fusion coefficients.

\clearpage\section{Verified models}\small\begin{longtable}{r p{0.84\textwidth}}\toprule Line & One verified realization\\\midrule\endhead

1 & $\operatorname{Rep}(D_{32})$\\[3pt]

10 & $\operatorname{TY}(C_{12})^{C_2}$, positive parity\\[3pt]

24 & $\mathrm{SU}(2)_6^{C_2}$, odd parity coherence\\[3pt]

26 & $\mathrm{SU}(2)_{10}$\\[3pt]

28 & $\mathrm{PSU}(2)_{20}$\\[3pt]

30 & $\mathrm{SO}(8)_2$\\[3pt]

37 & $\mathrm{SO}(15)_2$\\[3pt]

45 & $\operatorname{Rep}(D_{38})$\\[3pt]

46 & $\mathrm{PSU}(2)_{21}$\\[3pt]

48 & $\mathrm{Vec}_{C_{16}}^{\omega_1}$, equivariantized by $a\mapsto15a$\\[3pt]

57 & $V_4$ zest of $\mathrm{SO}(8)_2$, parameters $(0,0,0,2)$\\[3pt]

64 & $\operatorname{TY}(C_3\times C_2^2)^{C_2}$, positive parity\\[3pt]

68 & $\mathrm{SU}(2)_6^{C_2}$, trivial parity coherence\\[3pt]

72 & $\operatorname{TY}(C_{12})^{C_2}$, negative parity\\[3pt]

75 & $\operatorname{TY}(C_{15})^{C_2}$, negative parity\\[3pt]

103 & $C_2$ zest of $\mathrm{SU}(2)_{10}$\\[3pt]

104 & $\mathcal C(G_{32},\omega,V_4,1)$, nonsplit Jordan construction\\[3pt]

105 & $V_4$ zest of $\mathrm{SO}(8)_2$, parameters $(0,0,0,1)$\\[3pt]

112 & $\mathrm{Vec}_{C_{16}}^{C_2}$, action $a\mapsto7a$\\[3pt]

113 & $\operatorname{TY}(C_3\times C_2^2)^{C_2}$, negative parity\\[3pt]

146 & $\mathrm{Vec}_{C_{16}}^{\omega_1}$, equivariantized by $a\mapsto7a$\\[3pt]

150 & $V_4$ zest of $\mathrm{SO}(8)_2$, parameters $(2,0,0,3)$\\[3pt]

157 & $\operatorname{TY}(C_3^2)^{C_2}$, positive parity\\[3pt]

171 & $\operatorname{TY}(C_{10})$\\[3pt]

176 & $\operatorname{TY}(C_3^2)^{C_2}$, negative parity\\[3pt]

179 & Izumi $Q_8$ near-group category, equivariantized by $C_3$\\[3pt]

182 & $\mathrm{Vec}_{C_{11}}$\\[3pt]

\bottomrule\end{longtable}\normalsize
Here $D_m$ denotes a dihedral group of order $m$. Parameters, exact phases,
basis bijections and source hashes are in the individual model certificates;
the abbreviated names in the table are not the complete construction data.
\clearpage\section{Coherent model constructions}

\subsection{Established models and explicit coherent actions}
The model certificate in \texttt{known\_models/selected\_models.json} gives
25 unitary realizations, with an exact unital basis bijection and all
$11^3=1331$ structure constants checked for each.  It proves existence of
one model per ring, not uniqueness of associators.  The standard families
are pointed and representation categories, Tambara--Yamagami categories,
$\mathrm{SU}(2)_{10}$ and its indicated even counterparts,
$\mathrm{SO}(15)_2$, $\mathrm{SO}(8)_2$, and coherent equivariantizations
and associative zestings of these categories.  The complete formulas,
primary references, and scope of the finite searches are in
\texttt{known\_models/THEORY.md}.

The rank equations determine all members of the pointed, TY, $A_1$,
odd-metaplectic, full-even-metaplectic and half-even-metaplectic families
at this rank.  Prime rank prevents a nontrivial Deligne product.
For strict $C_2$ actions on pointed categories, the rank formula
$(|G|+3|G^{C_2}|)/2=11$ leaves $(|G|,|G^{C_2}|)=(19,1)$ or $(16,2)$.
The latter family is checked for all fourteen groups of order sixteen,
and gives rings $1$ and $112$.  The twisted pointed actions below are
checked separately; the strict scan is not mistaken for a classification
of all group-theoretical models.

For every TY equivariantization used, the certificate records a
nondegenerate symmetric bicharacter and an automorphism preserving it.
The automorphism therefore preserves the standard TY associator, giving
an actual tensor action.  The two choices of multiplication coherence
are the identity and the natural parity automorphism, equal to $-1$ on
the exceptional object and $+1$ on group objects.  Exact Fourier
projection computes the resulting fusion coefficients.  This covers
$10,37,64,72,75,113,157,176$.  In particular a based automorphism is
never used as an unverified categorical action.

\paragraph{A uniform zesting construction.}
In $\mathrm{SO}(8)_2$, the four neutral invertibles form a Tannakian
subcategory $\operatorname{Rep}(V_4)$.  Fix a coherent gauge in which
their associator and mutual braiding are trivial.  With respect to the
$V_4$ grading, every bilinear map
$\lambda:V_4\times V_4\longrightarrow V_4$ is an invertible-object
$2$-cocycle.  Its scalar associative-zesting obstruction is identically
one, since it involves only the displayed associator and mutual
braiding.  Thus
\[
 X\star Y=X\otimes Y\otimes\lambda(|X|,|Y|)
\]
has a unitary coherent associator.  The $4^4$ bilinear choices produce
four based rings, $30,57,105,150$.  The last is additional to the
single-$C_2$ trials.  The code checks all products, but the categorical
existence follows from this obstruction calculation, not merely from
associativity of the computed ring.  Ordinary $C_2$ zesting also supplies
ring $103$ from $\mathrm{SU}(2)_{10}$; the scalar obstruction vanishes
because $H^4(C_2,U(1))=0$.

\paragraph{The spin-current pair $24,68$.}
Let $\mathcal C=\mathrm{SU}(2)_6$ and let
\[
 \sigma(a)=\begin{cases}a&a\text{ even},\\6-a&a\text{ odd}.
 \end{cases}
\]
This dimension-preserving based involution is not accepted by itself.
Indeed the naive all-positive tensorator gives nonassociative proposed
fusion coefficients.  We use the standard real unitary quantum-Racah
associator to construct its tensorator explicitly.

Put $u=2\cos(\pi/8)$, so $u^4-4u^2+2=0$ and
$1847759/10^6<u<1847760/10^6$.  The quantum integers satisfy
$[1]=1$, $[2]=u$, $[j+1]=u[j]-[j-1]$.  The square of each quantum-Racah
$F$ entry lies in $\mathbb Q(u)$.  The verifier compares all transformed
squares exactly and determines every nonzero sign by rational interval
arithmetic in this fixed real embedding.  At this level the usual Racah
sum has no vanishing denominator; terms whose numerator contains
$[8]$ vanish.  The explicit formula and checks are included in the
source rather than delegated to numerical eigenvalue calculations.

There are $1664$ nonzero and $16$ zero admissible entries.  The nonzero
ratios define a consistent linear system of rank $65$ in $71$ normalized
$\{\pm1\}$ tensorator coefficients.  The recorded solution satisfies
every monoidal-functor equation.  The two recorded natural isomorphisms
$\mu:\sigma^2\Rightarrow\mathrm{Id}$ are
\[
 \mu_a=1,\qquad \mu_a=(-1)^a.
\]
Both satisfy $\delta\mu=J\,\sigma(J)$ and
$\mu_{\sigma(a)}=\mu_a$, including the group-associativity coherence.
They therefore define two unitary $C_2$ actions.  Orbit/stabilizer
projection computes their equivariantizations: on a fixed-output fusion
space the multiplicity is $(\dim+\operatorname{tr})/2$, with trace
computed from the actual tensorator and stabilizer eigenvalues.  The
first action yields ring $68$, and the second ring $24$, with every
coefficient matched.  The full tensorator and both coherences are
part of the portable certificate.

\paragraph{Twisted cyclic actions giving $48,146$.}
Let $n=16$, $t=1$, and write $c(a,b)=\lfloor(a+b)/n\rfloor$ for
residues in $\{0,\ldots,n-1\}$.  Give the pointed category on $C_n$
the associator
\[
 \omega_t(a,b,c)=\zeta_n^{\,t a c(b,c)}.
\]
For $u^2\equiv1\pmod n$, set $p(a)=ua\bmod n$ and
$v=(u^2-1)/n$.  A tensorator $F(a)\otimes F(b)\to F(a+b)$ is
\[
 J(a,b)=\zeta_n^{-tua\lfloor ub/n\rfloor}.
\]
The carry identity gives
$\delta J=\omega_t/(p^*\omega_t)$ and shows that the exponent of
$J(a,b)J(p(a),p(b))$ is $-tvab$ modulo $n$.  Define
\[
 \mu_a=\zeta_{2n}^{\,tv a^2+2ka}.
\]
Since $n$ is even, its coboundary is the composite tensorator.  The
remaining action equation is $\mu_{p(a)}=\mu_a$.  It holds with
$(u,v,k)=(15,14,0)$ and $(7,3,4)$, respectively.

These actions have fixed points $0$ and $8$, with $\mu_8=1$ and
$J(8,8)=-1$.  The lift of $8$ thus squares to the nontrivial stabilizer
character, giving a cyclic group of four invertibles.  The two
categories are identified with rings $48$ and $146$.  The verifier
checks, for each action, $16^4$ cocycle equations, $16^3$ tensorator
equations, $16^2$ square-coherence equations, and $16$ action-associativity
equations, then all $1331$ fusion coefficients.  Every phase is a root
of unity, so both constructions are unitary.

The relevant standard results are Tambara--Yamagami's existence theorem;
Gustafson--Rowell--Ruan, \emph{Metaplectic Categories, Gauging and
Property F}, Sections~3.1 and~3.3; Delaney--Galindo--Plavnik--Rowell--Zhang,
\emph{Braided zesting and its applications}; Burciu--Natale,
\emph{Fusion rules of equivariantizations of fusion categories}, corrected
version~2; and the Kirillov--Reshetikhin quantum-$6j$ formula.  The
simple-current autoequivalence also appears in Edie-Michell's
classification of autoequivalences.  None of these references is used
to bypass the explicit categorical action coherence established above.

\subsection{An explicit nonsplit group-theoretical model}
For ring $R11\_1\_104$ we construct an actual unitary group-theoretical
category.  Put $A=\mathbb F_2^3$, $z=(1,0,0)$ and
\[
 \sigma(a_0,a_1,a_2)=(a_0+a_1,a_1+a_2,a_2).
\]
Let $G$ be the order-$32$ group on pairs $(a,e)$, $0\leq e<4$, with
\[
 (a,e)(b,f)=\bigl(a+\sigma^e(b)+z\lfloor(e+f)/4\rfloor,
                         e+f\pmod4\bigr).
\]
Thus a lift $t$ of the cyclic generator satisfies $t^4=z$.  Let
$H=\{((0,a_1,a_2),0)\}\cong V_4$.

The file \texttt{known\_models/jordan\_gt\_model.json} specifies a
normalized cocycle $\omega:G^3\to\mu_8$ by its full exponent array.
All $32^4$ cocycle equations, all normalization equations, and
$\omega|_{H^3}=1$ are checked exactly.  Consequently the category of
$\mathbb C[H]$-bimodules in $\operatorname{Vec}_G^\omega$ is a unitary
fusion category $\mathcal C(G,\omega,H,1)$ of dimension $32$.
Its restriction to $A$ is the explicit cocycle
\[
 \omega(a,b,c)=(-1)^{a_0(b_0c_0+b_1c_1+b_2c_2+b_1c_2)}.
\]

The full relative tensor product is computed with the actual cocycle
phases, by projective characters of $H\cap gHg^{-1}$.  The derivation
from left and right bimodule actions is given in
\texttt{known\_models/PROOF\_JORDAN32.md}; the primary replay is
\texttt{known\_models/jordan\_verify.py} and an independent tensor
implementation is provided in \texttt{nonintegral/audit104}.
The model has eleven simples, dimensions $(1^4,2^7)$, and all
$1331$ fusion coefficients match $R11\_1\_104$ under the basis map
\[
 (0,4,8,5,9,1,2,3,10,6,7).
\]
Existence rests on the verified cocycle and bimodule construction,
not on associativity of a guessed fusion table.  This proves one
categorification, without asserting uniqueness of its associator.

% Embeddable research section. Requires amsmath, amssymb, amsthm, booktabs,
% array, and hyperref. Bibliography entries are in NONINTEGRAL-bib.tex.
\section{Dimension characters, central restrictions, and a quaternionic model}
\label{sec:r11-nonintegral}

Throughout this section, $R_n$ abbreviates the supplied based ring
\texttt{R11\_1\_}$n$. The integer $n$ is its physical line number in the
rank-eleven multiplicity-one input, and basis indices run from $0$ to $10$,
with $0$ the unit. All statements concern the entire based multiplication
tensor. In particular, an obstruction to a complex categorification is
distinguished from an obstruction that initially assumes a spherical or
pseudo-unitary structure.

The results below exclude thirty-eight rings as Grothendieck rings of
complex fusion categories. They also construct a unitary categorification
of $R_{179}$, of dimension $72$. The categorical ingredients are established
results; the uniform applications, explicit witnesses, and finite input
identifications are recorded here with exact replay certificates. These
results do not assert a classification of associators on any realizable
based ring.

\subsection{Exact passage from arbitrary categorifications to spherical ones}
\label{subsec:r11-signed-bridge}

We first explain the scope of the computational reduction used below.
Sphericalize a hypothetical complex fusion category $\mathcal C$, and
choose a lift of each simple object, with the unit and duality normalized.
Write $\varepsilon$ for the central sign object. In multiplicity one, the
quotient of the sphericalization's Grothendieck ring by
$[\varepsilon]+1$ has structure constants
\[
 T_{ij}^{k}=\sigma_{ij}^{k}N_{ij}^{k},\qquad
 \sigma_{ij}^{k}\in\{1,-1\}
 \quad\text{when }N_{ij}^{k}=1.
\]
These constants satisfy associativity, the unit identities, cyclic trace
symmetry, and reversal with duals. The trace triple $(i,j,k^*)$ is identified
only under cyclic rotation and
$(a,b,c)\mapsto(c^*,b^*,a^*)$; arbitrary permutations, or commutativity of
the signed quotient, are not imposed. Changing lifts replaces
$T_{ij}^{k}$ by $s_i s_j s_k T_{ij}^{k}$, where
$s_0=1$ and $s_i=s_{i^*}\in\{1,-1\}$.

\paragraph{Completeness of the sign calculation.}
Put $\sigma=(-1)^x$ on each oriented orbit. Every associativity identity
becomes an exact equation
\[
 \sum_m c_m(-1)^{\langle m,x\rangle}=0.
\]
If one coefficient has absolute value equal to the sum of the absolute
values of the others, the positive solution and equality in the triangle
inequality force the participating sign monomials to agree. Each resulting
linear relation over $\mathbb F_2$ is replayed from its full associativity
identity. An independently checked gauge matrix proves that the chosen
normalization meets every gauge orbit. All remaining assignments are then
enumerated and checked.

For the fifty-nine inputs in the recorded nonintegral audit, at most one
binary variable remains after this reduction. The complete domain has
seventy-one assignments: fifty-nine positive original tables and twelve
alternative associative signed tables. Every alternative is excluded by
its complete scalar-character inventory. Thus every categorically possible
signed table is gauge equivalent to the original positive table.

For completeness, an associative signed table with these reciprocity
identities is a finite-dimensional semisimple $*$-algebra in its regular
representation. The common kernel of its regular commutators is the sum
of the one-dimensional representation spaces. Proposed scalar characters
are checked directly on every product. Their distinct irreducible
separating polynomials, whose degrees sum to the dimension of this common
kernel, prove that no scalar character is omitted. A spherical dimension
character must be nonzero on every simple and equal on dual pairs. The
remaining exact obstructions use the algebraic integrality of
$\delta/f_E$ and the necessary containment
$f_E\in\mathbb Q(d_i)$; see
\cite[Theorem~2.13 and Corollary~2.15]{OstrikPivotal}.
For a signed quotient these conditions are applied in the actual
sphericalization: its dimension is $2\delta$, its two algebra summands
are the original and signed rings, and their formal codegrees are doubled.
Consequently the same field-containment conditions apply to both summands.

The signed conclusion has a categorical consequence, rather than merely
a numerical one. Choose the lifts for which every nonzero coefficient is
positive. Their products contain only chosen lifts, so they span a fusion
subcategory of the sphericalization. Forgetting the lifts gives a tensor
equivalence with $\mathcal C$. Its spherical structure therefore descends
to $\mathcal C$; compare \cite[Section~5.1]{OstrikPivotal}.

\paragraph{Original dimension characters.}
The same complete-inventory check is performed on the fifty-nine original
rings. Actual formal codegrees, rather than uncorrected regular Casimir
eigenvalues, are used: on a matrix block belonging to $E$, the latter
eigenvalue is $(\dim E)f_E$. The scalar part is extracted first and the
remaining block sizes are checked by the dimension of the semisimple
algebra. Exact norm and field-containment certificates then reject the
inadmissible dimension characters.

For $R_2,R_{51},R_{124},R_{147}$ no spherical dimension character remains.
For all surviving characters except those of
\[
 R_{36},R_{86},R_{95},R_{135},R_{136},R_{170},
\]
the inventory supplies a field embedding and a tensor sign character
making every simple dimension positive. Galois conjugation of the
algebraic categorical data and twisting the spherical pivotal structure
by that sign character then give positive dimensions, necessarily equal
to Frobenius--Perron dimensions. Thus a Galois conjugate is pseudo-unitary;
see \cite[Section~8.4]{ENO}. This assertion does not claim that the original
category is pseudo-unitary. The six exceptions have an additional mixed
spherical dimension character and will be treated directly.

The exact replay checks fifty-nine sign records, twelve alternative signed
character records, and fifty-nine original character records, with each
record bound to the full input tensor hash. It reconstructs the sign
deductions independently. The legacy sign checker also recomputes the
twelve small alternative character calculations; their results are then
checked separately by direct homomorphism identities and completeness,
without trusting that character-discovery calculation. This full replay,
not a stored Boolean conclusion, is a dependency of the general-complex
exclusions below.

\subsection{The spherical dimension lattice}
\label{subsec:r11-dimension-lattice}

Let $d_i$ be the simple dimensions of a spherical fusion category,
$\delta=\sum_i d_i d_{i^*}$ its global dimension, and
\[
 L_d=\sum_i\mathbb Z d_i
\]
the additive lattice generated by its simple dimensions. For an
irreducible representation $E$ of its based ring with formal codegree
$f_E$, the induction theorem
\cite[Theorem~2.13]{OstrikPivotal} supplies a simple
$A_E\in\mathcal Z(\mathcal C)$ such that
\begin{equation}\label{eq:r11-central-dimension}
 \dim A_E=\frac{\delta}{f_E},\qquad
 [F(A_E):\mathbf1]=\dim E,
\end{equation}
where $F$ is the forgetful functor. Since $F$ preserves dimension,
\begin{equation}\label{eq:r11-lattice-condition}
 \frac{\delta}{f_E}\in L_d
 \qquad\text{for every }E.
\end{equation}
This condition holds for every spherical dimension character, including
characters with negative dimensions or rational global dimension.

An exact certificate for failure of
\eqref{eq:r11-lattice-condition} is a rational linear functional
$\ell$ on the character field satisfying
\[
 \ell(d_i)\in\mathbb Z\quad\text{for every }i,
 \qquad \ell(\delta/f_E)\notin\mathbb Z.
\]
The verifier evaluates these expressions directly; it does not require
Hermite normal form or a lattice search. Each irreducible character-field
factor covers its entire Galois orbit, since both the dimension condition
and the multiset of formal codegrees are Galois invariant.

Here are all thirty-four inputs rejected by the complete spherical
inventory together with these lattice witnesses. Superscripts in the
dimension types denote multiplicities.
\begin{enumerate}
\item Type $1^4,2^3,(1+\sqrt5)^4$:
\[
 2,31,51,59,69,83,107,108,122,124,147,151.
\]
The Frobenius--Perron branch has
$D=40+8\sqrt5$ and $L=\mathbb Z+\mathbb Z\sqrt5$.
The actual formal codegree $40-8\sqrt5$ requires
\[
 \frac{D}{40-8\sqrt5}=\frac{3+\sqrt5}{2}\notin L.
\]
The four inputs $2,51,124,147$ were already excluded by their complete
spherical character inventories. For the other eight, all surviving
branches are covered by exact separating functionals.

\item Type $1^4,2^3,(2+2\sqrt2)^4$:
\[
 32,62,88,110,137,138,152,154.
\]
Here $D=64+32\sqrt2$ and $L=\mathbb Z+2\mathbb Z\sqrt2$.
The formal codegree $32$ requires
\[
 \frac{D}{32}=2+\sqrt2\notin L.
\]

\item Type $1^4,2,a^2,b^2,c^2$:
\[
 34,70,89,96,115,123,139,167.
\]
Let $a$ be the largest real root of $a^3-12a-8=0$, and put
$b=a^2/2-2$ and $c=a+2$.
The numbers $1,a,b$ are rationally independent, and
\[
 L=\mathbb Z+\mathbb Za+\mathbb Zb,
 \qquad D=48+12a+12b.
\]
The formal codegree $24$ requires
\[
 \frac{D}{24}=2+\frac{a+b}{2}\notin L.
\]

\item Type $1^4,2,(1+\sqrt5)^4,(3+\sqrt5)^2$:
\[
 36,86,95,135,136,170.
\]
The Frobenius--Perron branch has $D=60+20\sqrt5$ and
$L=\mathbb Z+\mathbb Z\sqrt5$; the formal codegree $40$ requires
$(3+\sqrt5)/2\notin L$.
For each of these six rings there is also a mixed spherical character
of global dimension $\delta=40$. Its simple dimensions still generate
$\mathbb Z+\mathbb Z\sqrt5$. The actual formal codegrees
$60\pm20\sqrt5$ now require
\[
 \frac{40}{60\pm20\sqrt5}=\frac{3\mp\sqrt5}{2}
 \notin\mathbb Z+\mathbb Z\sqrt5.
\]
Thus these branches are excluded directly, without an unjustified
pseudo-unitary reduction.
\end{enumerate}

The finite replay checks forty-four separating functionals over all
remaining character branches. Four of the thirty-four inputs have no
remaining spherical branch; the other thirty require at least one lattice
witness. Combined with Section~\ref{subsec:r11-signed-bridge}, these are
exclusions of arbitrary complex categorifications.

\subsection{A two-character central-restriction obstruction}
\label{subsec:r11-two-character}

The four further rings
\[
 R_3,\quad R_{52},\quad R_{142},\quad R_{149}
\]
have dimension type $1^4,2^3,x^4$, where
$x=2+2\sqrt2$, and global Frobenius--Perron dimension
$D=64+32\sqrt2$. For these four, the complete signed and character audit
reduces existence of a complex categorification to existence of a
pseudo-unitary one. We therefore work with positive simple dimensions.

Let $\chi$ be the Galois conjugate of the Frobenius--Perron character.
It has values $1$ on the invertibles, $2$ on the two-dimensional simples,
and $2-2\sqrt2$ on the four $x$-dimensional simples. Its formal codegree is
$f_\chi=64-32\sqrt2$. A second scalar character $\varphi$ has formal
codegree $f_\varphi=16$ and takes one of the values $\pm\sqrt2$ on each
$x$-dimensional simple. In input basis order, complete witnesses for this
second character are as follows, with $q=\sqrt2$:
\begin{center}
\small
\renewcommand{\arraystretch}{1.2}
\begin{tabular}{@{}rcl@{}}
\toprule
Input & $x$-dimensional indices & $(\varphi_0,\ldots,\varphi_{10})$\\
\midrule
$3$   & $4,5,6,7$ & $(1,1,1,0,-q,q,q,-q,0,1,-2)$\\
$52$  & $4,5,7,8$ & $(1,1,1,0,-q,q,-2,q,-q,0,1)$\\
$142$ & $1,2,3,4$ & $(1,q,q,-q,-q,-2,1,0,1,1,0)$\\
$149$ & $3,4,5,6$ & $(1,1,1,q,q,-q,-q,0,-2,0,1)$\\
\bottomrule
\end{tabular}
\end{center}
Every product identity and both formal codegrees are checked exactly.

Write
\[
 \beta_\chi=\frac{D}{f_\chi}=3+2\sqrt2=1+x,
 \qquad
 \beta_\varphi=\frac{D}{f_\varphi}=4+2\sqrt2=2+x.
\]
Since both characters are one-dimensional,
\eqref{eq:r11-central-dimension} says that the restrictions of their
associated central simples contain the unit exactly once. Comparing
coefficients of $\sqrt2$ and then rational parts, using positivity, forces
\begin{equation}\label{eq:r11-forced-restrictions}
 F(A_\chi)=\mathbf1+X,
 \qquad
 F(A_\varphi)=\mathbf1+g+Y,
\end{equation}
where $X,Y$ have dimension $x$ and $g$ is a nontrivial invertible.
Indeed the irrational coefficient permits exactly one $x$-dimensional
summand, and the remaining rational weight after the prescribed unit is
respectively zero or one.

The correspondence between based-ring characters and central simples,
together with symmetry of the center's $S$-matrix, gives
\begin{equation}\label{eq:r11-reciprocity}
 \beta_\chi\,\chi(F(A_\varphi))
 =\beta_\varphi\,\varphi(F(A_\chi));
\end{equation}
see \cite[Example~2.9 and the discussion preceding Theorem~2.13]
{OstrikPivotal}. But
\[
 \chi(F(A_\varphi))=4-2\sqrt2,
 \qquad
 (3+2\sqrt2)(4-2\sqrt2)=4+2\sqrt2.
\]
Thus \eqref{eq:r11-reciprocity} forces
$\varphi(\mathbf1+X)=1$, hence $\varphi(X)=0$. This contradicts the
displayed character values on every possible $X$. It excludes all four
rings by the same argument, without enumerating full induction matrices.
Together with the preceding thirty-four, the number of distinct
general-complex exclusions in this section is thirty-eight.

\subsection{A coherent quaternionic model for $R_{179}$}
\label{subsec:r11-quaternionic-model}

We now give an existence proof for an integral input. Let $\mathcal D$ be
one of Izumi's unitary near-group categories for the quaternion group
$Q_8=\{\pm1,\pm i,\pm j,\pm k\}$. Its fusion rules are
\begin{equation}\label{eq:r11-Q8-base}
 g\rho=\rho g=\rho,
 \qquad \rho^2=\sum_{g\in Q_8}g+2\rho.
\end{equation}
Thus $\operatorname{FPdim}(\rho)=4$ and
$\operatorname{FPdim}(\mathcal D)=24$.
Existence, with the explicit categorical data used below, follows from
\cite[Theorems~1.1 and~6.1, Lemma~6.13]{Izumi}.

\paragraph{The full integral categorical tuple.}
Use the faithful quaternionic representation on $K=\mathbb C^2$ given by
\[
 I=\begin{pmatrix}i&0\\0&-i\end{pmatrix},
 \qquad J=\begin{pmatrix}0&1\\-1&0\end{pmatrix},
 \qquad K_0=IJ,
\]
so that $\pi(i)=I$, $\pi(j)=J$, and $\pi(k)=K_0$.
Let $j_0=J\circ\text{complex conjugation}$. For any third root of unity
$\zeta$, Izumi's admissible tuple is
\begin{equation}\label{eq:r11-integral-tuple}
 (K,V,U_K,j_1,j_2,\chi,l)
 =(\mathbb C^2,\pi,\pi,j_0,-\zeta j_0,\chi,0),
 \qquad
 \pi(g)\pi(h)=\chi_h(g)\pi(h)\pi(g).
\end{equation}
This is the integral tuple of Lemma~6.13, not an application of an
automorphism theorem restricted to irrational near-group categories.

Put
\begin{equation}\label{eq:r11-order-three-W}
 W=-\frac{1+I+J+K_0}{2}.
\end{equation}
Direct Gaussian-rational matrix calculations give
\[
 W^*W=1,\qquad W^3=1,\qquad \operatorname{tr}(W)=-1,
 \qquad Wj_0=j_0W.
\]
Conjugation by $W$ cyclically permutes $I,J,K_0$. Hence, for the order-three
group automorphism $\phi:i\mapsto j\mapsto k\mapsto i$,
\[
 W\pi(g)W^*=\pi(\phi(g))\qquad(g\in Q_8).
\]
These equations preserve both representations $V$ and $U_K$ in
\eqref{eq:r11-integral-tuple}, both antiunitaries $j_1,j_2$, and the
commutator bicharacter $\chi$ after relabeling by $\phi$. Since $l=0$, its
compatibility is automatic. This verifies every component of the
equivalence-of-tuples condition in \cite[Definition~4.8]{Izumi}.

\paragraph{Categorical action and coherence.}
For clarity, existence of the action can be read directly from the
reconstruction, not merely from a permutation of fusion rules.
The generating isometries are
$S_g\in\operatorname{Hom}(g,\rho^2)$ and
$T(v)\in\operatorname{Hom}(\rho,\rho^2)$, $v\in K$.
Define the Cuntz-algebra automorphism
\begin{equation}\label{eq:r11-Cuntz-action}
 \gamma(S_g)=S_{\phi(g)},\qquad
 \gamma(T(v))=T(Wv).
\end{equation}
The reconstruction equations \cite[(3.10)--(3.12)]{Izumi} are covariant
under this assignment. In (3.10), $\sum_gS_g$ and the contraction
$\sum_a T_a j_1(T_a)$ are preserved. Equation (3.11) is covariant because
the pointed representations intertwine. Each term in (3.12) is covariant
because $W$ intertwines $V,U_K$, commutes with $j_1,j_2$, and preserves
$l=0$. Therefore $\gamma$ conjugates the pointed endomorphisms by $\phi$
and commutes with $\rho$. The reconstruction theorem
\cite[Theorem~4.9]{Izumi} supplies the resulting unitary tensor
autoequivalence of $\mathcal D$.

All morphisms are generated by these isometries and their adjoints
\cite[Remark~3.6]{Izumi}. Since $\phi^3=1$ and $W^3=1$, the cube of
\eqref{eq:r11-Cuntz-action} is the identity on the Cuntz generators,
hence on the reconstructed category. Identity compositors consequently
give a coherent, strict $C_3$ action. No use is made of Izumi's Section~13
automorphism results, which have a different, irrational scope.

\paragraph{Equivariant simples and products.}
The unitary equivariantization
$\mathcal C=\mathcal D^{C_3}$ has dimension $72$.
The fixed invertibles $1,-1$ each have three equivariant structures,
giving six invertibles $L_{(g,a)}$, $g\in\{1,-1\}$,
$a\in\mathbb Z/3$. The free orbits
$\{i,j,k\}$ and $\{-i,-j,-k\}$ give two objects $U,V$ of dimension three.
The fixed object $\rho$ has three lifts $P_a$ of dimension four.
The rank and dimension type are therefore $11$ and $1^6,3^2,4^3$.

The pointed part is computed by equivariant homogeneous fibers. Explicitly,
\begin{align*}
 L_{(g,a)}L_{(h,b)}&=L_{(gh,a+b)},\\
 L_{(1,a)}U&=U L_{(1,a)}=U,
 &L_{(1,a)}V&=V L_{(1,a)}=V,\\
 L_{(-1,a)}U&=U L_{(-1,a)}=V,
 &L_{(-1,a)}V&=V L_{(-1,a)}=U,\\
 U^2=V^2&=\sum_{a\in\mathbb Z/3}L_{(-1,a)}+U+V,
 &UV=VU&=\sum_{a\in\mathbb Z/3}L_{(1,a)}+U+V.
\end{align*}
The remaining products are
\begin{align}
 L_{(g,a)}P_b=P_bL_{(g,a)}&=P_{a+b},\label{eq:r11-fixed-rho}\\
 UP_b=P_bU=VP_b=P_bV&=P_0+P_1+P_2,\label{eq:r11-orbit-rho}\\
 P_aP_b&=L_{(1,a+b)}+L_{(-1,a+b)}+U+V\notag\\
 &\hspace{12mm}+P_{a+b+1}+P_{a+b+2}.
 \label{eq:r11-rho-square}
\end{align}
Indices are read modulo three. In \eqref{eq:r11-rho-square}, the fixed
pointed channels have phase one because $S_1,S_{-1}$ are fixed by
$\gamma$. The action on $\operatorname{Hom}(\rho,\rho^2)$ is $W$, whose
eigencharacters are the two nontrivial characters of $C_3$. An inverse
convention for the action replaces $W$ by $W^{-1}$ and leaves this
multiset unchanged. In the right absorption channels of
\eqref{eq:r11-fixed-rho}, no additional phase occurs: $\pi(\pm1)$ commutes
with $W$. Free three-cycles contribute each character once, yielding
\eqref{eq:r11-orbit-rho}. Finally $P_a^*=P_{-a}$, so exactly one
four-dimensional simple is self-dual.

In the model basis
\[
 L_{(1,0)},L_{(1,1)},L_{(1,2)},
 L_{(-1,0)},L_{(-1,1)},L_{(-1,2)},U,V,P_0,P_1,P_2,
\]
the model-to-input bijection is
\begin{equation}\label{eq:r11-179-bijection}
 (0,7,8,9,1,2,5,6,10,4,3).
\end{equation}
Every one of the $11^3=1331$ multiplication coefficients agrees under
this bijection with $R_{179}$. Its canonical tensor SHA-256 is
\begin{center}\small
\texttt{a9492a728c3c2e5d59d92d723a5349782a}\\[-1mm]
\texttt{962bc04ead4aacd0d7756e5d7c95f2}.
\end{center}
All three cube-root choices in \eqref{eq:r11-integral-tuple} provide
this based tensor. One such realization suffices for the existence claim.

\subsection{Finite verification and dependency order}
\label{subsec:r11-nonintegral-replay}

All following paths are relative to \texttt{work/rank11/nonintegral/}.
Run the four accepted entrypoints in this order, from the project root:
\begin{enumerate}
\item \nolinkurl{verify_reduction.py}: complete signed and scalar replay;
59 sign records, 12 alternative signed character records, and
59 original character records.
\item \nolinkurl{verify_spherical_lattice.py}: 34 excluded inputs,
including the four already rejected scalar inventories, and
44 exact separating functionals.
\item \nolinkurl{verify_two_character.py}: the four explicit character
pairs, formal codegrees, forced restrictions, and their contradictions.
Its rational arithmetic is implemented directly in $\mathbb Q(\sqrt2)$.
\item \nolinkurl{integral179/verify.py}: an independent reconstruction of
the equivariant fusion tensor by
\[
 \dim H^{C_3}=\frac13\sum_{t=0}^2\operatorname{Tr}(t\mid H)
\]
on every underlying near-group fusion space. The exceptional
two-dimensional space has character $(2,-1,-1)$; the other channels
are the one-dimensional absorption or $S_g$ spaces, with their orbit
permutations. This verifier does not import the model constructor.
It checks the exact action matrix, the source hash, all $1331$ tensor
coefficients, and all $11^4=14641$ associativity identities.
\end{enumerate}

The first replay imports the project's prior oriented-sign and exact
character verifiers. The lattice replay imports the shared exact-number
field helpers. The two-character and quaternionic model verifiers use
only the Python standard library. Critical conditions use explicit
exceptions and remain active under \texttt{python -O}. In particular,
the second and third checks alone are not certificates of a
general-complex exclusion: they require the full first replay, not merely
its summary flags. No timeout, failed model search, numerical tolerance,
or incomplete finite domain is used in any accepted conclusion here.

\section{A central dihedral obstruction in dimension 120}
\label{sec:rank11-d120}

\paragraph{Theorem.}
The based rings \texttt{R11\_1\_40} and \texttt{R11\_1\_82} admit no
categorification by a fusion category over $\mathbb C$.

\paragraph{Exact hypotheses.}
Both rings are commutative of rank eleven, with dimensions
$1^2,2,3^2,4^6$. Their nontrivial invertible $g$ has order two;
its multiplication permutation exchanges the two invertibles, exchanges
the two dimension-three simples, and fixes the other seven simples.
The basis elements of dimension less than four form a closed subring
of dimension twenty-four. No identification of this subring with a
particular category is assumed: in the second input its two dimension-three
simples are dual to each other, unlike the representation ring of $S_4$.

For a scalar character $\chi$, put
$f_\chi=\sum_i\chi(b_i)\chi(b_i^*)$ and $\beta_\chi=120/f_\chi$.
The eleven values $\beta_\chi$ are
\[
 1,2,2,8,8,8,8,8,15,30,30.
\]
Both characters with $\beta_\chi=2$ satisfy $\chi(g)=1$.
Finally the induction Gram coefficient $B_{gg}$ is eleven.
The accompanying exact verifier checks these statements directly from
the source tensor and a complete character table, including every
character multiplication identity. No restriction-frame search is used.

\paragraph{A forced central dihedral category.}
Suppose that $\mathcal C$ is a categorification. Its integrality supplies
the canonical positive spherical structure by
\cite[Proposition~8.24]{ENO}. Let $F:Z(\mathcal C)\to\mathcal C$ be
forgetful, with right adjoint $I$. By \cite[Theorems~2.5 and~2.13]{OstrikPivotal},
the simple constituents $A_\chi$ of $I(\mathbf1)$ occur once, have
$\dim(A_\chi)=\beta_\chi$, satisfy $[F(A_\chi):\mathbf1]=1$, and have
trivial twist. The character identification preceding Theorem~2.13 gives
\[
 S_{A_\chi,Y}=\beta_\chi\chi(F(Y))
 \qquad(Y\text{ simple in }Z(\mathcal C)),
\]
where $S$ is unnormalized. For the two dimension-two constituents $A,B$,
positivity forces $F(A)=F(B)=\mathbf1\oplus g$, whence
\[
 S_{A,A}=S_{A,B}=S_{B,B}=4.
\]
These equalities force trivial monodromy: the balancing formula is a
positive weighted sum of roots of unity, and equality with the sum of
the positive weights forces every eigenvalue to be one. Therefore the
subcategory $\mathcal T=\langle A,B\rangle$ is symmetric and has trivial
twist. It is Tannakian; see also \cite{DGNO}.

All objects of $\mathcal T$ forget to sums of $\mathbf1$ and $g$.
If a simple central $Y$ with this property contains $\mathbf1$ after
forgetting, it lies in $I(\mathbf1)$, with unit multiplicity one. Write
$F(Y)=\mathbf1\oplus m g$. The induction Gram equality gives
$m^2\le B_{gg}=11$, so $\dim(Y)\le4$. The displayed $\beta$ list
therefore leaves only $Y=\mathbf1,A,B$. In particular, the only simple
central object supported solely on $\mathbf1$ is $\mathbf1$ itself.

If instead $F(Y)=m g$, then $F(Y^*Y)=m^2\mathbf1$; all constituents of
$Y^*Y$ must be the central unit. Simplicity of $Y$ gives $m^2=1$.
Such a $Y$ is invertible and forgets to $g$. There is at most one such
lift, since the ratio of two lifts forgets to the unit; call it $z$ if
it exists. It must occur in $\mathcal T$: the dimension-four object
$A^*A$ contains the unit exactly once, and its remaining dimension three
cannot be filled using only the dimension-two simples $A,B$.
Thus $\mathcal T$ has exactly $\mathbf1,z,A,B$, of dimensions $1,1,2,2$.
It is consequently $\operatorname{Rep}(D_{10})$, where $D_{10}$ has order
ten. Its unique nontrivial invertible $z$ is a central boson with $F(z)=g$.

\paragraph{De-equivariantization and a faithful grading.}
The subcategory $\mathcal E=\langle z\rangle\simeq\operatorname{Rep}(C_2)$
embeds in $\mathcal C$ under $F$. By
\cite[Proposition~2.10]{ENOSolvable}, de-equivariantization gives a
category $\mathcal D$ of dimension sixty with
$\mathcal C\simeq\mathcal D^{C_2}$.
A free orbit $\{X,gX\}$ gives one simple of dimension $\dim(X)$;
a fixed $X$ gives two simples of dimension $\dim(X)/2$.
Here the fixed stabilizer algebra has two one-dimensional representations,
since $H^2(C_2,\mathbb C^\times)=0$, and its sign automorphism interchanges
them. The resulting dimension list is
\[
 \mathcal D:\quad 1^3,2^{12},3,
 \qquad \operatorname{rank}(\mathcal D)=16.
\]

Since $\mathcal T$ is symmetric, it lies in $\mathcal E'$.
The central equivalence $\mathcal E'_{C_2}\simeq Z(\mathcal D)$ sends
$\mathcal T_{C_2}$ to $\operatorname{Rep}(C_5)$.
Every object of this subcategory forgets to a multiple of the unit:
$\mathbf1,g$ become the unit under de-equivariantization, and the images
of $A,B$ become two copies of it. Their four nontrivial $C_5$-character
summands must each forget to the unit by positivity.
Proposition~2.9(i) of \cite{ENOSolvable} now supplies a faithful $C_5$
grading of $\mathcal D$, whose components all have dimension twelve.

The closed dimension-twenty-four subcategory above de-equivariantizes
to a subcategory $\mathcal K\subset\mathcal D$ of dimensions $1,1,1,3$.
Its invertibles form $C_3=\langle a\rangle$, and its unique other simple
$X$ is self-dual and satisfies
\[
 aX=Xa=X,\qquad X^2=\mathbf1+a+a^2+2X.
\]
Thus $\mathcal K$ has the fusion ring of $\operatorname{Rep}(A_4)$.
It equals its adjoint category, so every grading restricts trivially to
$\mathcal K$. Equality of dimensions gives $\mathcal D_e=\mathcal K$.

\paragraph{The Morita class of every possible neutral category.}
By \cite[Theorem~1.6]{ENOSolvable}, $\mathcal K$ is solvable. A solvable
tower consists of extensions and equivariantizations by cyclic groups of
prime order \cite[Definition~1.2 and Proposition~4.4]{ENOSolvable}.
Its final step cannot be an extension because $\mathcal K$ has trivial
universal grading. Hence $\mathcal K=\mathcal L^{C_p}$ for a prime $p$.
The embedded $\operatorname{Rep}(C_p)$ forces $p=3$, and
$\operatorname{FPdim}(\mathcal L)=4$.

The simple equivariant dimensions are $[C_3:H]\dim(\pi)\dim(Y)$.
All projective irreducible representations of the cyclic stabilizers $H$
are one-dimensional. Since the only dimensions in $\mathcal K$ are one
and three, a fixed $Y$ has dimension one or three, the latter impossible
in a category of dimension four. A nonfixed $Y$ has an orbit of length
three and likewise has dimension one. Thus
$\mathcal L=\operatorname{Vec}_H^\eta$ for a group $H$ of order four.
The action on $H$ is nontrivial, since otherwise equivariantization would
have twelve invertibles. Thus $H=V_4$, with $C_3$ cycling its nonidentity
elements.

The pointed crossed product $\mathcal L\rtimes C_3$ is Morita equivalent
to $\mathcal K$ by \cite[Proposition~2.2]{ENOSolvable}. Its group of
invertibles is the extension of $V_4$ by $C_3$ with this nontrivial action;
the coprime extension splits, giving $A_4$. Consequently
\[
 \mathcal K\text{ is Morita equivalent to }\operatorname{Vec}_{A_4}^\Omega
\]
for some three-cocycle $\Omega$, with no triviality assumption on $\Omega$.

\paragraph{No order five in the twisted Brauer--Picard group.}
Let $\mathcal P=\operatorname{Vec}_{A_4}^\Omega$.
The twisted classification \cite[Theorem~1.2, or Theorem~4.12]{NaiduNikshych}
parametrizes Lagrangian subcategories of $Z(\mathcal P)$ by normal abelian
subgroups and invariant alternating $\Omega$-bicharacters.
The only possible subgroups are $1,V_4$. The first contributes one choice;
for $V_4$, the ratio of two admissible twisted bicharacters is an ordinary
alternating bicharacter, of which there are exactly two. Hence there are
at most three Lagrangian subcategories.

The stabilizer of the canonical $\operatorname{Rep}(A_4)$ preserves its
regular algebra $\operatorname{Fun}(A_4)$: a connected etale algebra of
dimension twelve in $\operatorname{Rep}(A_4)$ corresponds uniquely to the
free transitive $A_4$-set. Conversely the regular algebra generates that
subcategory. Theorem~3.3 and Remark~3.4 of \cite{MarshallNikshych} identify
this stabilizer with the induction image $I(A_4,\Omega)$.
The elementary bound from tensor autoequivalences already suffices:
their kernel over the underlying automorphism group is
$H^2(A_4,\mathbb C^\times)=C_2$ and their image is contained in
$\operatorname{Aut}(A_4)=S_4$, so the stabilizer order divides $48$.

More precisely, Proposition~4.5 of \cite{MarshallNikshych}, valid for
arbitrary $\Omega$, gives
\[
 Z(A_4)\longrightarrow H^2(A_4,\mathbb C^\times)
 \longrightarrow I(A_4,\Omega)
 \longrightarrow\operatorname{OutStab}(\Omega)\longrightarrow1.
\]
Since $Z(A_4)=1$ and $\operatorname{Out}(A_4)=C_2$, the stabilizer order
is two or four. These group-theoretic data are recorded in
\cite[Section~8.3]{NikshychRiepel}; its untwisted Brauer--Picard computation
is not being applied to a twisted category.
Orbit--stabilizer, with orbit size at most three, shows that
$\operatorname{BrPic}(\mathcal P)\simeq
\operatorname{Aut}_{\mathrm{br}}(Z(\mathcal P))$ has no element of order
five. Morita invariance gives the same statement for $\mathcal K$.

\paragraph{Contradiction.}
A faithful graded extension determines a homomorphism
\[
 C_5\longrightarrow\operatorname{BrPic}(\mathcal K),
 \qquad h\longmapsto[\mathcal D_h].
\]
It is trivial, so every component is equivalent to the regular
$\mathcal K$-bimodule category and has rank four. Thus
$\operatorname{rank}(\mathcal D)=20$, contrary to the forced rank sixteen.
This proves the theorem for both inputs.\hfill$\square$

\section{A cyclic extension restriction obstruction}
\label{sec:rank11-c4}

\paragraph{Theorem.}
The based ring \texttt{R11\_1\_153} admits no categorification over
$\mathbb C$. The argument applies to every complex fusion category,
without a braiding, pivotality, or unitarity assumption.

\paragraph{A general restriction criterion.}
Let $\mathcal D$ be a fusion category, $G$ a finite group, $K\le G$, and
$c:G\to\operatorname{BrPic}(\mathcal D)$ a fixed homomorphism.
Put $A=\operatorname{Inv}Z(\mathcal D)$ with the action induced by $c$.
Suppose a reference $G$-extension realizing $c$ exists and
\[
 \operatorname{res}:H^2(G,A)\longrightarrow H^2(K,A)
 \quad\text{is zero.}
\]
Then the $K$-subextensions of all $G$-extensions realizing $c$ have
isomorphic based fusion rings, respecting their gradings and neutral
subcategories up to the chosen identifications.

Indeed, systems of multiplication functors form an $H^2(G,A)$-torsor
by \cite[Theorem~8.5]{ENOHomotopy}. Comparing two existing extensions,
their product systems differ by an $A$-valued $2$-cocycle. Restricting
the multiplication functors restricts this cocycle. If its restricted
class is zero, componentwise bimodule equivalences remove the
difference, so the restricted fusion coefficients agree. The remaining
$H^3(G,\mathbb C^*)$-freedom changes associators by scalars and leaves
multiplication functors unchanged
\cite[Theorem~8.9]{ENOHomotopy}. Duality is determined by the coefficient
of the unit, so it is preserved. This argument compares existing
extensions and requires no separate obstruction-vanishing assertion.

\paragraph{Exact structural hypotheses.}
Both \texttt{R11\_1\_104} and \texttt{R11\_1\_153} have dimension type
$(1^4,2^7)$, universal grading $C_4$, and homogeneous ranks $(5,2,2,2)$.
Their neutral ring is the Tambara--Yamagami ring on $V_4$; write $m$ for
its dimension-two simple. In a generating component there is a simple
$X$ with
\[
 XX^*=\mathbf1+h_L+m,\qquad X^*X=\mathbf1+h_R+m,
 \qquad h_L\ne h_R\ne\mathbf1,\quad h_L\ne\mathbf1.
\]
The source verifier derives the adjoint subring and all homogeneous
components directly from the tensors and checks every homogeneous
product, rather than only products of chosen representatives.

\paragraph{The neutral category is necessarily $\operatorname{Rep}(D_8)$.}
Every categorification of the neutral ring is
$\operatorname{TY}(V_4,\chi,\tau)$, with $\chi$ a nondegenerate symmetric
bicharacter and $\tau=\pm\tfrac12$.
The canonical internal-end algebras $XX^*$ and $X^*X$ have the form
$\mathbf1+h+m$. Their multiplication maps $h\otimes m\to m$ and
$m\otimes h\to m$ are nonzero: the $h$-summand arises from an invertible
stabilizing action on $X$. Associativity on $(h,m,h)$ therefore forces
$\chi(h,h)=1$. There are two distinct nonzero isotropic elements,
$h_L,h_R$. A nonalternating nondegenerate symmetric form on
$\mathbb F_2^2$ has exactly one, so $\chi$ is alternating.
The two remaining TY categories are $\operatorname{Rep}(D_8)$ and
$\operatorname{Rep}(Q_8)$, with groups of order eight.
But $\operatorname{BrPic}(\operatorname{Rep}(Q_8))\cong S_3$
\cite[Remark~9.4]{NikshychRiepel}. Any homomorphism from $C_4$ to this
group kills its degree-two element, forcing that homogeneous bimodule
to be regular of rank five, contrary to its rank two.
Thus $\mathcal D\simeq\operatorname{Rep}(D_8)$. This includes all neutral
associators and both TY signs.

\paragraph{The Brauer--Picard homomorphism is unique up to neutral equivalence.}
By \cite[Section~9.1]{NikshychRiepel},
$\operatorname{BrPic}(\mathcal D)\cong S_4$.
Reindexing the alternating TY data realizes the six isometries of $V_4$
as tensor autoequivalences. Their map into the Brauer--Picard group is
injective: its kernel consists of inner tensor autoequivalences, which
are tensor isomorphic to the identity in a braided category.
This $S_3$ is the stabilizer of the regular left-module class, a point
$p$ of the four-point $S_4$ action.
All three nonidentity components have rank two, so the first, second,
and third powers of the grading generator move $p$. Its image is
therefore a $4$-cycle. The $S_3$ fixing $p$ conjugates the six $4$-cycles
transitively: its intersection with the centralizer of one is trivial.
Consequently every relevant $c$ can be aligned with the reference by
an actual tensor autoequivalence of $\mathcal D$.

\paragraph{The coefficient module.}
Here $A=\operatorname{Inv}Z(\mathcal D)\cong C_2^3$ and, as an $S_4$-module,
\[
 A\cong\mathbb F_2^4/\langle(1,1,1,1)\rangle.
\]
For completeness this action can be identified intrinsically. Let $z$
be the nontrivial central invertible with underlying tensor unit: its
half-braiding acts as $+1$ on the four invertibles of $\mathcal D$ and
$-1$ on $m$. Every tensor autoequivalence fixes $z$, so its $S_4$ orbit
has size one or four. If its orbit had size one, the $C_4$ action of
the explicit \texttt{R11\_1\_104} model would fix $z$.
Relative-center equivariantization \cite{CentersGraded} would then lift
$z$ to a central invertible with underlying unit: the projective
linearization obstruction lies in $H^2(C_4,\mathbb C^*)=0$.
But tensor automorphisms of the identity of that category are
characters of its universal grading $C_4$, hence restrict trivially
to its adjoint category $\mathcal D$. This contradicts the nontrivial
restriction $z$.
Thus the orbit consists of four distinct nonzero vectors. They span
$A$, since a two-dimensional space over $\mathbb F_2$ has only three
nonzero vectors. The unique relation is $S_4$-invariant and is their
sum zero, proving the displayed module identification.

\paragraph{Vanishing of the restriction.}
Let $t$ be a $4$-cycle and $N=t-1$ on $A$. This is a single unipotent
Jordan block of size three:
\[
 N^3=0,\qquad \operatorname{rank}N=2,\qquad
 \operatorname{rank}N^2=1,\qquad \ker N=\operatorname{im}N^2.
\]
The $C_4$ norm is $1+t+t^2+t^3=N^3=0$, so cyclic cohomology gives
\[
 H^2(C_4,A)=\ker N,\qquad
 H^2(\langle t^2\rangle,A)=\ker N^2/\operatorname{im}N^2.
\]
Even-degree restriction sends an invariant representative to itself.
Since $\ker N=\operatorname{im}N^2$, this restriction is zero.
This computation is also verified without invoking the periodic
resolution: an independent normalized bar-cochain program enumerates
all $128$ normalized $C_4$ $2$-cocycles, finds the $64$ coboundaries,
and supplies a $C_2$ $1$-cochain trivializing the restriction of every
cocycle. It also verifies the six conjugacy representatives above.

\paragraph{Reference model and contradiction.}
The reference \texttt{R11\_1\_104} is the explicit group-theoretical
category $\mathcal C(G,\omega,H,1)$ provided by the finite certificate
\texttt{jordan\_gt\_model.json}. Here $|G|=32$, $H\cong V_4$, and
$\omega$ is a normalized $\mu_8$-valued $3$-cocycle trivial on $H$.
The group is the nonsplit Jordan extension of $C_2^3$ by $C_4$,
with $t^4$ equal to the fixed nonzero vector.
The certificate includes the complete group table and cocycle,
all simple labels, the full model fusion tensor, and a unit-preserving
basis bijection to the original census tensor. All $32^4$ cocycle
identities and all fusion coefficients are checked. An independent
implementation reconstructs all $1331$ coefficients from free algebra
bimodules, balanced tensor quotients, and graded intertwiner systems;
it imports no projective-character convolution routine. This audit
passes and is executed by the exclusion verifier.

In \texttt{R11\_1\_104}, the degree-two simples have source labels
$9,10$ and are both self-dual. In \texttt{R11\_1\_153}, those labels
are $1,2$ and they are dual to each other. Thus the respective even
subrings have seven and five self-dual simples and are not isomorphic
as based rings. A hypothetical categorification of the latter ring
has the forced neutral category and the same $4$-cycle homomorphism
as the reference, up to neutral equivalence. The restriction criterion
forces its even subring to equal that of the reference, a contradiction.

\paragraph{Scope and reproducibility.}
The exclusion uses the new general restriction criterion and a positive
reference model, not a claim that a bounded model search is exhaustive.
The earlier failed cocycle ans\"atze have no role in the proof.
The source-bound records are \texttt{c4\_reduction.json} and
\texttt{c4\_exclusion.json}; the independent normalized-cochain audit is
\texttt{c4\_torsor\_independent.json}. The categorical neutral-module
audit is \texttt{C4\_NEUTRAL\_MODULE\_AUDIT.md}.

\section{Complete status index}
\paragraph{Excluded (155).} 2, 3, 4, 5, 6, 7, 8, 9, 11, 12, 13, 14, 15, 16, 17, 18, 19, 20, 21, 22, 23, 25, 27, 29, 31, 32, 33, 34, 35, 36, 38, 39, 40, 41, 42, 43, 44, 47, 49, 50, 51, 52, 53, 54, 55, 56, 58, 59, 60, 61, 62, 63, 65, 66, 67, 69, 70, 71, 73, 74, 76, 77, 78, 79, 80, 81, 82, 83, 84, 85, 86, 87, 88, 89, 90, 91, 92, 93, 94, 95, 96, 97, 98, 99, 100, 101, 102, 106, 107, 108, 109, 110, 111, 114, 115, 116, 117, 118, 119, 120, 121, 122, 123, 124, 125, 126, 127, 128, 129, 130, 131, 132, 133, 134, 135, 136, 137, 138, 139, 140, 141, 142, 143, 144, 145, 147, 148, 149, 151, 152, 153, 154, 155, 156, 158, 159, 160, 161, 162, 163, 164, 165, 166, 167, 168, 169, 170, 172, 173, 174, 175, 177, 178, 180, 181.

\paragraph{Unitary models (27).} 1, 10, 24, 26, 28, 30, 37, 45, 46, 48, 57, 64, 68, 72, 75, 103, 104, 105, 112, 113, 146, 150, 157, 171, 176, 179, 182.

Every number above denotes the suffix in \texttt{R11\_1\_i}. The original
input has uncompressed SHA-256
\begin{quote}\small\ttfamily
3186c410888e0bc7f30c9e423e0184a156\linebreak
18246bc9f4d0ef503e708c4081eff8
\end{quote}

\begin{thebibliography}{99}
\bibitem{ENO}
P.~Etingof, D.~Nikshych, and V.~Ostrik,
\emph{On fusion categories},
Annals of Mathematics \textbf{162} (2005), 581--642.
\href{https://arxiv.org/abs/math/0203060}{arXiv:math/0203060}.


\bibitem{OstrikPivotal}
V.~Ostrik,
\emph{Pivotal fusion categories of rank 3},
Moscow Mathematical Journal \textbf{15} (2015), 373--396.
\href{https://arxiv.org/abs/1309.4822}{arXiv:1309.4822}.


\bibitem{Izumi}
M.~Izumi,
\emph{A Cuntz algebra approach to the classification of near-group categories},
arXiv:1512.04288v2 (2015).
\href{https://arxiv.org/abs/1512.04288}{arXiv:1512.04288}.


\bibitem{ENOSolvable}
P.~Etingof, D.~Nikshych and V.~Ostrik,
\emph{Weakly group-theoretical and solvable fusion categories},
Adv. Math. \textbf{226} (2011), 176--205;
arXiv:0809.3031v2.


\bibitem{DGNO}
V.~Drinfeld, S.~Gelaki, D.~Nikshych and V.~Ostrik,
\emph{On braided fusion categories I},
Selecta Math. \textbf{16} (2010), 1--119;
arXiv:0906.0620v3.


\bibitem{NaiduNikshych}
D.~Naidu and D.~Nikshych,
\emph{Lagrangian subcategories and braided tensor equivalences of twisted
quantum doubles of finite groups},
Comm. Math. Phys. \textbf{279} (2008), 845--872;
arXiv:0705.0665v2.


\bibitem{MarshallNikshych}
I.~Marshall and D.~Nikshych,
\emph{On the Brauer--Picard groups of fusion categories};
arXiv:1603.04318v3.


\bibitem{NikshychRiepel}
D.~Nikshych and B.~Riepel,
\emph{Categorical Lagrangian Grassmannians and Brauer--Picard groups of
pointed fusion categories},
J. Algebra \textbf{411} (2014), 191--214;
arXiv:1309.5026v2.


\bibitem{ENOHomotopy}
P.~Etingof, D.~Nikshych and V.~Ostrik,
\emph{Fusion categories and homotopy theory},
Quantum Topol. \textbf{1} (2010), 209--273;
arXiv:0909.3140v2.


\bibitem{CentersGraded}
S.~Gelaki, D.~Naidu and D.~Nikshych,
\emph{Centers of graded fusion categories},
Algebra Number Theory \textbf{3} (2009), 959--990;
arXiv:0905.3117.

% NikshychRiepel is provided by D120-bib.tex in the combined report.
\end{thebibliography}
\end{document}